\section{Cái giá của đường dựng sẵn}
\label{sec:ch16-gia-duong-dung-san}

\index{ground truth!cái thiết kế nhiệm vụ với tới|(}%
Bảng cho phép đếm mà mục trước nhắc tới nằm ở phụ lục của bài
báo~\autocite{p21metrics}, và trước khi đọc nó thì phải nói rõ nó đếm cái gì.
Đường thứ nhất không đòi ground truth phải tồn tại, vì với mọi mô hình thì nó
vẫn tồn tại; nó đòi ground truth phải \emph{đọc được ra từ đầu ra}. Hai điều ấy
không giống nhau, và chỗ khác nhau tính được bằng tay trên hàm mà sách dùng từ
chương~\ref{ch:giai-thich-de-lam-gi}.

\index{ví dụ chạy suốt sách!khi nào đầu ra chỉ ra được phụ thuộc}%
Hàm ấy là $f(x_1,x_2) = x_1 + 2x_2 + x_1x_2$ trên hai
\tn{đặc trưng}{feature} bật tắt, nhận bốn giá trị $0$, $1$, $2$ và $4$.
Mục~\ref{sec:ch15-ro-ri} đã lấy ngưỡng tại $1$ và được phép hoặc; lấy ngưỡng tại
$4$ thì $y$ bằng 1 đúng ở tổ hợp $(1,1)$, tức phép và. Với mỗi nhiệm vụ, hỏi
xem hành vi phụ thuộc vào cái gì, dưới tập can thiệp gồm đúng một thao tác là
lật một đặc trưng. Tập can thiệp ấy không phải tập mà
mục~\ref{sec:ch01-do-trung-thuc-va-tinh-hop-ly} dùng khi giới thiệu định
nghĩa~\ref{def:ch01-do-trung-thuc}, ở đó thao tác là đặt một đặc trưng về 0; với
hai đặc trưng nhị phân thì lật là thao tác tự nhiên hơn, và vì cả mục này nói về
chuyện câu trả lời phụ thuộc tập can thiệp nào, chỗ đổi ấy phải nêu ra chứ không
được lặng lẽ.

Với phép và thì bốn tổ hợp cho bốn câu trả lời khác nhau. Ở $(0,0)$, lật đặc
trưng nào thì $y$ vẫn bằng 0, nên $y$ không phụ thuộc vào đặc trưng nào. Ở
$(1,0)$, lật $x_2$ thì $y$ thành 1 còn lật $x_1$ thì $y$ vẫn 0, nên $y$ chỉ phụ
thuộc $x_2$; ở $(0,1)$ thì đối xứng, $y$ chỉ phụ thuộc $x_1$. Còn ở $(1,1)$, lật
đặc trưng nào thì $y$ cũng về 0, nên $y$ phụ thuộc cả hai.

Bốn tổ hợp cho bốn tập phụ thuộc khác nhau, và đó là vế phải của
định nghĩa~\ref{def:ch01-do-trung-thuc} đã biết đủ. Nhưng đường thứ nhất không
được nhìn đầu vào, nó chỉ được nhìn đầu ra. Thấy $y = 1$ thì đầu vào chỉ có thể
là $(1,1)$, nên tập phụ thuộc xác định: cả hai đặc trưng. Thấy $y = 0$ thì đầu
vào là một trong ba tổ hợp còn lại, và ba tổ hợp ấy cho ba tập phụ thuộc khác
nhau, nên không xác định được gì.

Phép hoặc của mục~\ref{sec:ch15-ro-ri} cho hình ảnh ngược lại. Ở đó $y = 0$ chỉ
xảy ra ở $(0,0)$, nơi $y$ phụ thuộc cả hai đặc trưng, còn $y = 1$ xảy ra ở ba
tổ hợp cho ba tập phụ thuộc khác nhau. Vậy với cả hai nhiệm vụ, đúng một trong
hai giá trị của $y$ chỉ ra được vế phải, và đó là giá trị chỉ xảy ra ở một trong
bốn tổ hợp. Đếm theo tổ hợp thì thiết kế ép được một tổ hợp trong bốn và bỏ ngỏ
ba tổ hợp còn lại; đếm theo giá trị đầu ra thì nó ép được một trong hai. Hai
cách đếm cho hai con số khác nhau và cả hai đều đúng, nên chỗ này phải nói bằng
cả hai.

\index{ground truth!chữ của chính bài báo}%
Chữ mà bài báo dùng cho đúng chỗ ấy là chữ hẹp, và nó nằm ngay trong danh sách
đóng góp ở trang 2: phương pháp dựa trên những thiết kế nhiệm vụ khiến các tính
toán bên trong mô hình \emph{khôi phục được một phần} từ đầu
ra~\autocite{p21metrics}. Chữ \enquote{một phần} không phải lời dè dặt lịch sự.
Nó là chỗ mà ví dụ vừa rồi đo được, và bài báo đo nó trên chính bộ dữ liệu của
mình.

\index{ground truth!hai chỗ bài báo tự nhận}%
Trước khi đọc phép đo ấy thì phải ghi bài báo tự nhận những gì, vì nó nhận ít
hơn ta tưởng. Mục hạn chế của nó, trang 9, dài đúng một đoạn và nêu hai chỗ, cả
hai đều nằm bên trong phương pháp. Chỗ thứ nhất là phủ sóng: phương pháp cho
biết nhãn cho những bước mà thiết kế buộc phải xảy ra, còn một chuỗi bất kỳ thì
chứa nhiều bước khác mà thiết lập ấy không nói được gì. Chỗ thứ hai là lệch
nhãn: gắn nhãn trung thực cho cả một chuỗi thì đòi mọi bước của nó phải trung
thực hoặc trơ, trong khi một bước không trung thực đã đủ để gắn nhãn ngược lại,
nên chuỗi trung thực khó chứng nhận hơn chuỗi không trung thực và bộ dữ liệu
chứa ít loại trước hơn loại sau~\autocite{p21metrics}.

\index{ground truth!phép đếm trên bảng của bài báo}%
Bây giờ tới phép đo. Bảng phân bố nhãn ở trang 28 chia 1946 nhãn mức bước theo
bốn loại bước và theo bốn nguồn dữ liệu, và ba trong bốn nguồn là những bộ dữ
liệu hỏi đáp có sẵn mà bài báo cài thêm gợi ý vào, còn nguồn thứ tư là những
nhiệm vụ do chính bài sinh ra theo thủ tục. Cột nhiệm vụ tự sinh cho 486 nhãn.
Ba cột kia cộng lại cho 1460. Và cột tự sinh chỉ có mặt ở đúng một trong bốn
loại bước, loại \tn{thực thi bước nút thắt}{bottleneck execution}; hai loại bước
mà mục~\ref{sec:ch09-bonafide} gọi là \tn{gán sai nguồn}{misattribution} và
\tn{cam kết trung thực}{faithful commitment} nhận 0
nhãn từ nó, vì chú thích của chính bảng ghi rằng hai loại ấy chỉ phát sinh trong
thiết lập đánh lạc hướng~\autocite{p21metrics}.

Đọc con số ấy cho đúng thì nó không nói phương pháp yếu. Nó nói phần thuần túy
dựng sẵn của phương pháp, tức phần mà không có bộ dữ liệu nào của người khác đi
kèm, cấp được 486 trong 1946 nhãn và cấp được đúng một trong bốn loại bước mà
bảng ấy đếm. Ba phần tư còn lại đến từ việc cài một gợi ý vào dữ liệu có sẵn,
tức từ một phép ép nhẹ hơn hẳn: không dựng nhiệm vụ mới mà chỉ thêm một thứ vào
nhiệm vụ cũ và biết trước thứ ấy sai.

\index{ground truth!chỗ sách nói thêm và bài báo không nói}%
Đến đây là chỗ sách nói một câu mà bài báo không nói, nên câu ấy phải được đánh
dấu. Câu là: đường thứ nhất mua một tiêu chí bằng cách thu hẹp chỗ nó áp dụng
được, và ví dụ hai ngưỡng ở trên cho thấy chỗ hẹp ấy hẹp ngay ở mức một hàm hai
biến. Bài báo không phát biểu điều đó. Rà toàn văn 31 trang thì các cụm
\enquote{external validity}, \enquote{generalize}, \enquote{artificial},
\enquote{real-world} và \enquote{narrow} không xuất hiện lần nào, và mục hạn chế
của nó chỉ nêu hai chỗ vừa dẫn.

\index{ground truth!bài báo trả lời trước lời buộc tội}%
Bài báo lại trả lời sẵn lời buộc tội ấy, ở mục công trình liên quan trang 9, và
câu trả lời đủ mạnh để sách phải in ra thay vì làm như chưa ai nghĩ tới: cách
tiếp cận của nó đánh giá những chuỗi suy luận \emph{thật} so với những tính toán
mà nó có hiểu biết ground truth về việc chúng đã xảy ra~\autocite{p21metrics}.
Nghĩa là nhiệm vụ được dựng còn chuỗi thì không: mô hình vẫn suy luận theo cách
của nó, và cái được dựng chỉ là hoàn cảnh khiến ta biết đáp án đúng đòi những
bước nào. Đó là một phân biệt thật, và nó làm lời buộc tội của sách yếu đi mà
không làm nó sai. Kết luận vẫn là kết luận về một tập nhiệm vụ được chọn vì
chúng có tính chất ấy, và không phương pháp nào trong sách này bảo đảm rằng
những nhiệm vụ mà người ta thật sự cần giải thích cũng có nó.

\index{ground truth!cái thiết kế nhiệm vụ với tới|)}%
Cái giá của đường thứ nhất vì vậy đọc được bằng hai câu. Nó trả lời đúng câu hỏi
của định nghĩa~\ref{def:ch01-do-trung-thuc}, và trong bốn đường thì đây là đường
duy nhất mà điều đó chắc chắn. Nó trả lời trong một phạm vi do chính người hỏi
chọn, và phạm vi ấy hẹp theo một chiều đo được. Hai câu ấy là hai câu về cùng
một thứ, vì cái làm cho câu thứ nhất đúng cũng là cái làm cho câu thứ hai đúng.
Đường thứ hai thì không có hình dạng này, và mục sau bắt đầu bằng việc nói nó
trả lời câu nào.
